\documentclass[11pt,a4paper]{article}

% Packages
\usepackage[utf8]{inputenc}
\usepackage[T1]{fontenc}
\usepackage{amsmath,amsfonts,amssymb,amsthm}
\usepackage{mathtools}
\usepackage{graphicx}
\usepackage{hyperref}
\usepackage{cleveref}
\usepackage{booktabs}
\usepackage{algorithm}
\usepackage{algpseudocode}
\usepackage{subcaption}
\usepackage[margin=1in]{geometry}

% Author affiliations
\usepackage{authblk}

% Bibliography
\bibliographystyle{plain}

% Theorem environments
\newtheorem{theorem}{Theorem}[section]
\newtheorem{proposition}[theorem]{Proposition}
\newtheorem{lemma}[theorem]{Lemma}
\newtheorem{corollary}[theorem]{Corollary}
\newtheorem{definition}[theorem]{Definition}
\newtheorem{remark}[theorem]{Remark}
\newtheorem{example}[theorem]{Example}

% Custom commands
\newcommand{\R}{\mathbb{R}}
\newcommand{\N}{\mathbb{N}}
\newcommand{\norm}[1]{\left\|#1\right\|}
\newcommand{\inner}[2]{\left\langle #1, #2 \right\rangle}
\newcommand{\pd}[2]{\frac{\partial #1}{\partial #2}}

% Title information
\title{%
  [Draft] Extensions to the Deep Fourier Residual Method\\
  \large A Follow-up Study
}

\author[1]{Author Name}
\author[2]{Collaborator Name}
\affil[1]{Institution 1}
\affil[2]{Institution 2}

\date{\today}

\begin{document}

\maketitle

\begin{abstract}
This paper extends the Deep Fourier Residual (DFR) method introduced by Taylor, Pardo, and Muga~\cite{taylor2023deep} for solving partial differential equations using neural networks. We explore [describe your follow-up idea here]. Our contributions include [list main contributions]. Numerical experiments demonstrate [summarize key results].
\end{abstract}

\section{Introduction}
\label{sec:introduction}
% Introduction section

Physics-Informed Neural Networks (PINNs) have emerged as a powerful paradigm for solving partial differential equations (PDEs) using deep learning~\cite{raissi2019physics}. The key idea is to embed physical constraints directly into the neural network training process through the loss function.

Recently, Taylor, Pardo, and Muga~\cite{taylor2023deep} introduced the Deep Fourier Residual (DFR) method, which uses the $H^{-1}$ dual norm of the weak residual as the loss function. This approach provides several advantages over traditional PINNs:
\begin{itemize}
    \item The loss is equivalent to the $H^1$ error for well-posed problems
    \item Reliable error estimation during training
    \item Applicability to problems with limited regularity
\end{itemize}

\subsection{Motivation}

% Describe the motivation for your follow-up work
% What limitations of DFR are you addressing?
% What new applications are you enabling?

[Describe the motivation for your extension here]

\subsection{Contributions}

The main contributions of this work are:
\begin{enumerate}
    \item [First contribution]
    \item [Second contribution]
    \item [Third contribution]
</enumerate>

\subsection{Organization}

The remainder of this paper is organized as follows. Section~\ref{sec:background} reviews the necessary background on PINNs and the DFR method. Section~\ref{sec:methodology} presents our proposed extensions. Section~\ref{sec:theory} provides theoretical analysis. Section~\ref{sec:experiments} describes our numerical experiments, and Section~\ref{sec:results} discusses the results. Finally, Section~\ref{sec:conclusion} concludes the paper.


\section{Background}
\label{sec:background}
% Background section

\subsection{Problem Setting}

We consider the abstract problem of finding $u \in U$ such that
\begin{equation}
    \mathcal{R}(u) = 0,
\end{equation}
where $\mathcal{R}: U \to V^*$ is the PDE operator mapping from a trial space $U$ to the dual of a test space $V$.

For example, Poisson's equation with homogeneous Dirichlet boundary conditions:
\begin{equation}
    -\Delta u = f \quad \text{in } \Omega, \qquad u = 0 \quad \text{on } \partial\Omega,
\end{equation}
can be expressed in weak form as: find $u \in H^1_0(\Omega)$ such that
\begin{equation}
    \inner{\nabla u}{\nabla v}_{L^2} = \inner{f}{v}_{L^2} \quad \forall v \in H^1_0(\Omega).
\end{equation}

\subsection{Physics-Informed Neural Networks}

Physics-Informed Neural Networks (PINNs)~\cite{raissi2019physics} approximate the solution using a neural network $u_\theta$ and minimize a loss function based on the PDE residual:
\begin{equation}
    \mathcal{L}_{\text{PINN}}(\theta) = \frac{1}{N} \sum_{i=1}^{N} |\mathcal{R}^s(u_\theta)(x_i)|^2 + \lambda \cdot \text{BC penalty},
\end{equation}
where $\mathcal{R}^s$ denotes the strong-form residual and $\{x_i\}$ are collocation points.

\subsection{Deep Fourier Residual Method}

The DFR method~\cite{taylor2023deep} uses the weak formulation and minimizes the $H^{-1}$ norm of the residual:
\begin{equation}
    \mathcal{L}_{\text{DFR}}(\theta) = \norm{\mathcal{R}^w(u_\theta)}_{H^{-1}}^2.
\end{equation}

For rectangular domains, the $H^{-1}$ norm can be computed efficiently using the Discrete Sine/Cosine Transform:
\begin{equation}
    \norm{R}_{H^{-1}}^2 = \sum_{k} \frac{|\hat{R}_k|^2}{1 + |k|^2},
\end{equation}
where $\hat{R}_k$ are the Fourier coefficients.

\subsubsection{Key Properties}

The DFR method enjoys the following properties:
\begin{proposition}[Error equivalence~\cite{taylor2023deep}]
    For well-posed problems, there exist constants $\gamma, M > 0$ such that
    \begin{equation}
        \frac{1}{M}\norm{\mathcal{R}(u)}_{V^*} \leq \norm{u - u^*}_{U} \leq \frac{1}{\gamma}\norm{\mathcal{R}(u)}_{V^*}.
    \end{equation}
\end{proposition}

This error equivalence ensures that minimizing the DFR loss is equivalent to minimizing the error, providing reliable training dynamics.

\subsection{Limitations of Current Approaches}

While the DFR method addresses several limitations of PINNs, some challenges remain:
\begin{itemize}
    \item Restriction to rectangular domains
    \item Extension to time-dependent problems
    \item Scalability to high dimensions
    \item [Add other limitations relevant to your work]
\end{itemize}

[Continue with background relevant to your specific extension]


\section{Methodology}
\label{sec:methodology}
% Methodology section

In this section, we present our proposed extensions to the Deep Fourier Residual method.

\subsection{Overview}

[Provide a high-level overview of your proposed method]

\subsection{Proposed Method}

[Describe your method in detail]

\subsubsection{Mathematical Formulation}

[Present the mathematical formulation]

Let $\mathcal{N}_\theta: \R^d \to \R$ denote the neural network with parameters $\theta$. Our proposed loss function is:
\begin{equation}
    \mathcal{L}(\theta) = \text{[Your loss function here]}
\end{equation}

\subsubsection{Algorithm}

\begin{algorithm}
\caption{[Your Algorithm Name]}
\label{alg:method}
\begin{algorithmic}[1]
\Require Neural network $\mathcal{N}_\theta$, problem data
\Ensure Trained parameters $\theta^*$
\State Initialize parameters $\theta$
\For{$k = 1, \ldots, K$}
    \State [Algorithm steps]
    \State Compute loss $\mathcal{L}(\theta)$
    \State Update $\theta \gets \theta - \eta \nabla_\theta \mathcal{L}$
\EndFor
\State \Return $\theta^*$
\end{algorithmic}
\end{algorithm}

\subsection{Implementation Details}

[Describe implementation details]

\subsubsection{Network Architecture}

We use a fully-connected feedforward neural network with:
\begin{itemize}
    \item Input dimension: $d$ (spatial dimension)
    \item Hidden layers: [specify]
    \item Activation function: [specify]
    \item Output dimension: 1 (scalar solution)
\end{itemize}

\subsubsection{Boundary Condition Enforcement}

For homogeneous Dirichlet conditions, we use the cutoff approach:
\begin{equation}
    u_\theta(x) = \phi(x) \cdot \tilde{u}_\theta(x),
\end{equation}
where $\phi(x) = 0$ on $\partial\Omega$ and $\tilde{u}_\theta$ is the raw network output.

\subsubsection{Training Procedure}

[Describe training procedure, optimizer choice, learning rate schedule, etc.]


\section{Theoretical Analysis}
\label{sec:theory}
% Theory section

In this section, we provide theoretical analysis of our proposed method.

\subsection{Main Results}

[State your main theoretical results]

\begin{theorem}[Main result]
\label{thm:main}
[State your main theorem here]
\end{theorem}

\begin{proof}
[Proof or sketch of proof. Full proof can go in appendix]
\end{proof}

\subsection{Error Analysis}

[Provide error analysis]

\begin{proposition}[Error bound]
\label{prop:error}
Under assumptions [A1]-[A3], the error satisfies:
\begin{equation}
    \norm{u_\theta - u^*}_{H^1} \leq C \cdot \text{[bound]}
\end{equation}
\end{proposition}

\subsection{Convergence Analysis}

[Analyze convergence properties]

\begin{theorem}[Convergence]
\label{thm:convergence}
[State convergence result]
\end{theorem}

\subsection{Comparison with DFR}

[Compare theoretical properties with the original DFR method]

\begin{remark}
[Discuss relationship to original DFR method]
\end{remark}


\section{Numerical Experiments}
\label{sec:experiments}
% Experiments section

We evaluate our proposed method on a series of benchmark problems, comparing against PINNs and the original DFR method.

\subsection{Experimental Setup}

\subsubsection{Problems}

We consider the following test problems:

\paragraph{Problem 1: [Name]}
[Describe first test problem]
\begin{equation}
    -\Delta u = f \quad \text{in } \Omega
\end{equation}
with exact solution $u^*(x) = \text{[exact solution]}$.

\paragraph{Problem 2: [Name]}
[Describe second test problem]

\paragraph{Problem 3: [Name]}
[Describe third test problem]

\subsubsection{Implementation}

All experiments were implemented using:
\begin{itemize}
    \item Python 3.10
    \item TensorFlow/JAX/PyTorch (specify which)
    \item Network architecture: [specify]
    \item Training: Adam optimizer, learning rate $\eta = 10^{-3}$, [epochs] iterations
\end{itemize}

Code is available at: \url{https://github.com/mathmode/pinns-dfr}

\subsubsection{Metrics}

We evaluate performance using:
\begin{itemize}
    \item $L^2$ error: $\norm{u_\theta - u^*}_{L^2}$
    \item $H^1$ error: $\norm{u_\theta - u^*}_{H^1}$
    \item Training loss
    \item Computational time
\end{itemize}

\subsection{Results}

[Present results, reference figures and tables]

\begin{table}[htbp]
\centering
\caption{Comparison of methods on benchmark problems}
\label{tab:results}
\begin{tabular}{lccc}
\toprule
Problem & PINNs & DFR & Ours \\
\midrule
Problem 1 & & & \\
Problem 2 & & & \\
Problem 3 & & & \\
\bottomrule
\end{tabular}
\end{table}

\begin{figure}[htbp]
\centering
% \includegraphics[width=0.8\textwidth]{figures/results.pdf}
\caption{[Figure caption]}
\label{fig:results}
\end{figure}

\subsection{Ablation Studies}

[Present ablation studies if applicable]

\subsection{Computational Cost}

[Discuss computational cost comparison]


\section{Results and Discussion}
\label{sec:results}
% Results and Discussion section

\subsection{Summary of Findings}

Our experiments demonstrate that [summarize main findings]:

\begin{enumerate}
    \item [Finding 1]
    \item [Finding 2]
    \item [Finding 3]
\end{enumerate}

\subsection{Analysis}

[Provide detailed analysis of results]

\subsubsection{Comparison with PINNs}

[Compare with PINNs]

\subsubsection{Comparison with DFR}

[Compare with DFR]

\subsection{Discussion}

[Discuss implications of results]

\subsubsection{Advantages}

Our proposed method offers the following advantages:
\begin{itemize}
    \item [Advantage 1]
    \item [Advantage 2]
    \item [Advantage 3]
\end{itemize}

\subsubsection{Limitations}

We acknowledge the following limitations:
\begin{itemize}
    \item [Limitation 1]
    \item [Limitation 2]
\end{itemize}

\subsection{Practical Recommendations}

Based on our findings, we recommend:
\begin{itemize}
    \item [Recommendation 1]
    \item [Recommendation 2]
\end{itemize}


\section{Conclusion}
\label{sec:conclusion}
% Conclusion section

We have presented [brief summary of contributions].

\subsection{Summary}

The main contributions of this work are:
\begin{enumerate}
    \item [Contribution 1]
    \item [Contribution 2]
    \item [Contribution 3]
\end{enumerate}

Our numerical experiments on [problems] demonstrate that [key finding].

\subsection{Future Work}

Several directions for future work remain:

\begin{itemize}
    \item \textbf{Extension to general domains}: [describe]
    \item \textbf{Time-dependent problems}: [describe]
    \item \textbf{Higher dimensions}: [describe]
    \item \textbf{Nonlinear problems}: [describe]
\end{itemize}

\subsection{Reproducibility}

All code and data used in this paper are available at:
\begin{center}
    \url{https://github.com/mathmode/pinns-dfr}
\end{center}

The experiments can be reproduced using the provided configuration files and scripts. See the repository README for detailed instructions.


\section*{Acknowledgments}
[Acknowledgments here]

\bibliography{references}

\appendix
\section{Additional Proofs}
\label{app:proofs}
% Additional proofs and derivations

\section{Implementation Details}
\label{app:implementation}
% Code availability and implementation notes

\end{document}
